\usepackage{graphicx} 
\usepackage{verbatim}
\usepackage{hyperref}
\hypersetup{
    colorlinks,%
    citecolor=black,%
    filecolor=black,%
    linkcolor=black,%
    urlcolor=blue
}

\graphicspath{{../figures/}}

\mode<article>{\usepackage{fullpage}}
\mode<presentation>{\usetheme{default}}

\newcommand{\ftitle}[1]{\frametitle<presentation>{#1}}
\newcommand{\cmd}[1]{\begin{quote}{\tt #1}\end{quote}}
\newcommand{\lit}[1]{\texttt{#1}}

\title{UNIX Workshop 2010\\\url{http://uws.assembla.me}}
\author{
Angad~Singh \and 
Aldrian~Obaja \and
Dang~Cao~Khoa \and
Chong~Perry \and
Saiaswin~Saikannan \and
Jialong~Liu\\ \and
Vu~An~Hoa \and
Gary~Kwong \and
Melvin~Zhang
}

\begin{document}

\maketitle

\mode<presentation>{
  \begin{frame}
  \titlepage
  \end{frame}
}

\begin{frame}<presentation>
\begin{center}
\includegraphics[width=\textwidth]{jurassic_park}

\url{http://www.youtube.com/watch?v=dFUlAQZB9Ng}
\end{center}
\end{frame}

\begin{frame}
\ftitle{Jurassic Park (1993)}
\begin{quote}
``It's a UNIX system! I know this.''
\end{quote}
\begin{flushright}
-- Alexis ``Lex'' Murphy, Jurassic Park (1993)
\end{flushright}
\end{frame}

\section{Why UNIX?} %{{{
You may be running Windows on your desktop/laptop, but the next time you visit
Google or Facebook, the webpage you see is being served from a computer running
UNIX. According to \url{w3techs.com}, in August 2010, at least 67\% of all web
servers are running on UNIX.

UNIX is the name of the operating system developed by a group of AT\&T
researchers at Bell Labs in 1969.  Most modern operating systems, such as Linux,
Solaris, Mac OS X, and BSD (see Figure \ref{fig:modern}), descended from UNIX.
All of these operating systems provide a UNIX-like environment to the user. In
the rest of these notes, the term UNIX will refer to these modern day
descendants of the original AT\&T UNIX.   

\begin{frame}
\ftitle{Modern UNIX}
\begin{figure}
\begin{center}
\includegraphics[width=250pt]{unix_logos}
\caption{Modern UNIX-like operating systems (Linux, BSD, Solaris)}
\label{fig:modern}
\end{center}
\end{figure}
\end{frame}

\begin{frame}<presentation>{What do these sites have in common?}
\includegraphics[width=\textwidth]{site_logos} 
\pause

\begin{Large}
67\% of all web servers are running on UNIX\footnote{\url{http://w3techs.com}, August 2010}
\end{Large}
\end{frame}

\begin{frame}<presentation>{C was invented to write UNIX}
\begin{center}
\includegraphics[width=0.6\textwidth]{c}
\end{center}
\end{frame}

The story of UNIX is also tied in with the creation of the C programming
language.  Dennis Ritchie invented C in order to rewrite UNIX. Coincidentally,
starting from the 2010 intake, the School of Computing changed the programming
language used in its programming methodology module from Java to C. For those of
you doing Scheme, you will learn C when you start on your Operating System
module (CS2106).   

\begin{frame}<presentation>{You will be programming in UNIX}
\begin{center}
\includegraphics[width=0.8\textwidth]{programming}
\end{center}

\begin{itemize}
\item CS1010 labs - developing C programs in UNIX.  
\item CS1020 labs - developing Java programs in UNIX
\end{itemize}
\end{frame}

Finally, in addition to learning about programming methodology in the
introductory modules, you will be required to develop your C/Java programs in
the UNIX environment.

The basic workflow for writing programs is the edit-compile-run-test cycle (see
Figure \ref{fig:flow}).  For interpreted languages, such as Scheme, there is no
compile step.  In this workshop, we will be focusing on the parts of UNIX
related to the edit-compile-run-test cycle.  

\begin{frame}
\begin{figure}
\begin{center}
\includegraphics[page=1,width=250pt]{flow}
\caption{Workflow for writing programs}
\label{fig:flow}
\end{center}
\end{figure}
\end{frame}
%}}}

\section{Getting started} %{{{ 
SoC operates a number of UNIX servers where students can login and interactively
develop programs in the UNIX environment.  We will be accessing SoC's \lit{sunfire}
server from the PCs in the programming lab.  Use your NUSNET account to login to
the PC in the programming lab.  

% ACTIVITY
\begin{frame}{Activity: Login to NUSNET}
\begin{enumerate}
\item Press \lit{Ctrl-Alt-Delete}.
\mode<presentation>{
\begin{figure}
\begin{center}
\includegraphics[scale=0.3]{ctrl_alt_delete}
\end{center}
\end{figure}
}
\item Type in your NUSNET user name, password and select the NUSSTU domain.  
\mode<presentation>{
\begin{figure}
\begin{center}
\includegraphics[scale=0.3]{stu_login}
\end{center}
\end{figure}
}
\item Click on the \lit{Ok} button.  
\end{enumerate}
\end{frame}

Before you can login to \lit{sunfire}, you will need to create your UNIX account.
In the following activity, you will create your UNIX account.  

% ACTIVITY
\begin{frame}{Activity: Creating your UNIX account}
\mode<presentation>{
\begin{Large}
\url{https://mysoc.nus.edu.sg/~newacct}
\end{Large}

Login using your NUSNET user name and password.  
}
\mode<article>{
\begin{enumerate}
\item Login to \url{https://mysoc.nus.edu.sg/~newacct} using your NUSNET
user name and password.  
\item Decide your UNIX user name.  
\item Type in your new password (twice).
\item Submit your application.
\end{enumerate}
}
\end{frame}


SoC's \lit{sunfire} server runs on the Solaris operating system developed by Sun
Microsystem.  

%}}}

\section{Command line interface} %{{{

\begin{frame}
\ftitle{\lit{sunfire} server in the old Machine Room}
\begin{figure}
\begin{center}
\includegraphics[scale=0.5]{machine_room}
\caption{\lit{sunfire} server located in the Machine Room with our Central
Facilities staff. Clockwise from top-left: Tan Chee Sin, Tan Kwang Pon, Budiman
Tsjin and Lai Zit Seng (Systems programmer, ITU).}
\label{fig:sunfire}
\end{center}
\end{figure}
\end{frame}

In the next activity, we will connect to the \lit{sunfire} server (see Figure
\ref{fig:sunfire}) and access the UNIX system via the Secure Shell protocol
(ssh).

\begin{frame}{Activity: Connecting to \lit{sunfire}}
\begin{enumerate}
\item From the desktop, launch the \lit{SSH Secure Shell Client} application. 
\item Click on \lit{Quick Connect}\\ 
Host Name: \lit{sunfire.comp.nus.edu.sg}\\ 
User Name: your UNIX user name
\mode<presentation>{
\begin{figure}
\begin{center}
\includegraphics[scale=0.25]{ssh_sunfire}
\end{center}
\end{figure}
}
\item Click on \lit{Connect}.
\item Click on ``Yes'' at the Host identification dialog.  
\item Enter your UNIX password in the password dialog.  
\end{enumerate}
\end{frame}


% Linux/Mac OS X already comes with a UNIX-like environemnt,
% recommend windows users use Cygwin  

% Act: create a folder to put files for UWS

% Basic concepts:
% Launching program by typing name follow by parameters, programs as functions, f(x) ?
% Silence is golden: there is usually no output when a program runs successfully

% Activity:
% home directory
% ls
% pwd
% mkdir UNIXWorkshopFiles
% cd UN (tab completion)
% pwd
% cd ..
% cd UN (tab completion)


% Intro to Linux Comand Line, http://www.youtube.com/watch?v=6kxK8Fqcs5o

\begin{frame}
\ftitle{Command line interface}
\begin{figure}
\begin{center}
\includegraphics[width=250pt]{ssh_command_line}
\caption{Command line interface on \lit{sunfire}}
\label{fig:cli}
\end{center}
\end{figure}
\end{frame}

Once you have logged in, you will see a command line interface (see Figure
\ref{fig:cli}). 

Our next activity involves working with files and directories. On UNIX, files
are arranged in a hierarchy (see Figure \ref{fig:dir}) of directories.  The top
of the hierarchy is traditionally called root (written as a slash /).

\begin{frame}
\ftitle{UNIX Directory Tree}
\begin{figure}
\begin{center}
\includegraphics[scale=0.4]{file_home}
\end{center}
\caption{A subset of the UNIX directory tree showing home directories}
\label{fig:dir}
\end{figure}
\end{frame}

In this interface, we interact with the system by typing in a single line of
command, followed by the \lit{Enter} key.  A command consists of two parts: the
name of the program you wish to execute, followed by a list of arguments or
input to that program.  

\begin{frame}[fragile]
\ftitle{Parts of a command}
\begin{Large}
\begin{verbatim}
$ program argument1 argument2 ...
\end{verbatim}
\end{Large}
\end{frame}

Such an interface might seem rather primitive compared to the rich graphical
user interfaces (GUI) we are used to.  However, there are a number of advantages
to using the CLI. It is a reasonable interface for programs that do not need to
interact with the user and it makes it easy to invoke any program even if there
are thousands of programs installed. Once you are comfortable with this
interface, it can be much faster than using the GUI.  



In the following activity, we will use the command line to create a directory to
store files related to this workshop.  There is a concept of a \emph{working
directory}. Commands that operate on files or directories assumes that you are
modifying your working directory.  

% ACTIVITY
\begin{frame}[fragile]{Activity: Working with files and directories}
\begin{enumerate}
\item After login, you are place in your home directory, e.g. \lit{/home/m/melvin}
\item You can check your working directory using the \lit{pwd} command \cmd{pwd}
\item The \lit{ls} command shows you the files in your working directory \cmd{ls}
\item Now create a new directory called \lit{UNIXWorkshopFiles} \cmd{mkdir UNIXWorkshopFiles}
\item Switch to the new folder using the \lit{cd} command
\cmd{cd UNIXWorkshopFiles}
\item Use the \lit{pwd} again command to check your working directory \cmd{pwd}
\end{enumerate}
\end{frame}

%}}}

\section{Text editing} %{{{

\begin{frame}<presentation>
\includegraphics[page=2,width=\textwidth]{flow}
\end{frame}

There are many types of files in UNIX, the include:
\begin{itemize}
\item a document (report, essay etc.)
\item the text of a program written in some high-level programming language
\item instructions comprehensible directly to the machine and incomprehensible 
to a casual user, for example, a collection of binary digits (an executable or 
binary file)
\item a directory, containing information about its contents, which may be a 
mixture of other directories (subdirectories) and ordinary files.
\end{itemize}

By far the most common are text files, which are used extensively on a UNIX
system for storing system data, program configuration files, scripts and source
code. They are preferable to binary files because they can be easily
read/modified.  The most important program for working with text files is a
\emph{text editor}, a program to interactively create/edit text files.  Examples
of editors on UNIX are nano, Vim, and Emacs.

\begin{frame}
\ftitle{Text files are ubiquitous on UNIX}
\mode<presentation>{
Program source code are stored as text files. A good text editor can
dramatically improve your productivity.

\begin{figure}
\begin{center}
\includegraphics[width=0.75\textwidth]{nano}
\end{center}
\caption{Screenshot of nano}
\end{figure}
}
\end{frame}

In the following activity, we will use nano, \url{http://www.nano-editor.org/},
to edit a simple C program. Nano is a basic text editor that is easy to learn,
if you are interested in a more sophisticated text editor, you should try
learning Vim or Emacs.  

% ACTIVITY
\begin{frame}[fragile]{Activity: Text editing with nano}
\begin{enumerate}
\item Download the sample \lit{GCD.c} program from the UWS website using wget
\cmd{wget http://uws.assembla.me/GCD.c}
\item Edit the file using name \lit{nano}
\cmd{nano GCD.c}
\item Type in your name and matric number as indicated in the file.  
\item Save the file and exit nano by pressing \cmd{Ctrl-x}
\item Check the contents of the file using the \lit{cat} command
\cmd{cat GCD.c}
\end{enumerate}
\end{frame}

\begin{frame}<presentation>
\includegraphics[page=3,width=\textwidth]{flow}
\end{frame}

\begin{frame}{Activity: Compiling and running}
\begin{enumerate}
\item C programs are compiled using the \lit{gcc} compiler.
\cmd{gcc GCD.c}
\item To run a program, you must add \lit{./} in front of its name. The default
name used by gcc is \lit{a.out}
\item Run the GCD program
\cmd{./a.out}
\item Type in a pair of integers followed by the \lit{Enter} key, for example
\cmd{58 24}
\item Repeat step 4 as many times as you like.  
\item To force the program to end, press \lit{Ctrl-c}
\end{enumerate}


\end{frame}

%}}}

\section{Working with text files} %{{{

\begin{frame}<presentation>
\begin{center}
\includegraphics[page=4,width=\textwidth]{flow}
\end{center}
\end{frame}

\begin{frame}{Testing}
\begin{center}
\includegraphics[width=250pt]{bbtesting}
\end{center}
\end{frame}

\begin{frame}[fragile]{Testing on UNIX}
\begin{Large}
\begin{verbatim}
$ program < input > output
\end{verbatim}
\end{Large}
\end{frame}

% ACTIVITY
\begin{frame}[fragile]{Activity: Creating the input}
\begin{enumerate}
\item Instead of typing the input by hand as in the previous activity, we create
an input file using nano. \cmd{nano input}
\item Type in pairs of integers, one pair per line, for example
\begin{verbatim}
3 10
15 25
200 420
\end{verbatim}
\item Save and exit nano. 
\end{enumerate}
\end{frame}

\begin{frame}[fragile]{Activity: Creating the correct output}
\begin{enumerate}
\item Create an file for the correct output
\cmd{nano answer}
\item Type in the correct GCD for each pair of integers in the input
\begin{verbatim}
GCD(3,10)=1
GCD(15,25)=5
GCD(200,420)=20
\end{verbatim}
\item Save and exit nano.  
\item Run the GCD program on the test case. \cmd{./a.out < input > output}
\end{enumerate}
\end{frame}

\begin{frame}{When things go wrong}
\begin{center}
\includegraphics[width=250pt]{bug}
\end{center}
\end{frame}

\begin{frame}{Activity: Find the bugs!}

\begin{enumerate}
\item Open \lit{output} in nano and verify if it matches the correct answer.
\cmd{nano output}
\item If they differ, you've found the bug!, else try again with different
\lit{input/answer} files
\item Hint: manually checking whether two files are identical is boring, try
using the \lit{diff} command
\cmd{diff answer output}
\end{enumerate}
\end{frame}

%}}}

\section{Combing multiple tools} %{{{
UNIX is more than just a family operating systems. It also an approach towards
software that emphasizes creating small sharp tools that work well together.
This is embodied in the UNIX philosophy.   

\begin{frame}{The UNIX Philosophy}
\begin{quote}
Write programs that do one thing and do it well.

Write programs to work together.

Write programs to handle text streams, because that is a universal interface.
\end{quote}
\begin{flushright}
\mode<presentation>{\includegraphics[height=0.3\textheight]{mcilroy}\\}
-- Douglas McIlroy\\(inventor of UNIX pipes)
\end{flushright}
\end{frame}

The next activity demonstrates the UNIX philosophy, by making use
of three simple UNIX utility programs to analyse SMS messages.  

% ACTIVITY: SMS Word Count
\begin{frame}{Activity: SMS Word Count}
\mode<presentation>{
\begin{figure}
\begin{center}
\includegraphics[scale=0.5]{sms}
\end{center}
\end{figure}
}
Your friend from FASS is studying SMS language as part of a course project. She
collected a number of SMS messages and would like to find out the frequency of
each word. 
\end{frame}

\begin{frame}[fragile]
\ftitle{Activity: SMS Word Count}
\begin{columns}
\begin{column}{0.5\textwidth}
For example, given the following text file:
\begin{verbatim}
U wan 2 haf lunch i'm in da 
canteen now.
Haf u found him? I feel so 
stupid da v cam was working.
Where r we meeting?
I went to ur hon lab but no 
one is there.
\end{verbatim}
\end{column}
\begin{column}{0.5\textwidth}
The desired output is:
\begin{verbatim}
      .
      .
      .
      1 we
      1 went
      1 Where
      1 working.
      2 da
      2 I
\end{verbatim}
\end{column}
\end{columns}
\end{frame}

% Hugo: Input file is too large, maybe reduce to 5 lines of SMS
% Hugo: introduce grep
% introduce man to find the option of sort and uniq to use

\begin{frame}[fragile]
\frametitle{Activity: \lit{sort} and \lit{uniq}}
Two UNIX utility programs are related to our task.  

\begin{block}{\lit{sort}}
\begin{tabular}{c c c}
Input: & & Output: \\
dog    & & bat \\
bat    &$\longrightarrow$& cat \\
log    & & dog \\
cat    & & log \\
\end{tabular}
\end{block}

\begin{block}{\lit{uniq}}
\begin{tabular}{c c c}
Input: & & Output: \\
dog    & & dog \\
dog    &$\longrightarrow$& cat \\
cat    & & dog \\
cat    & & cat \\
dog    & & \\
cat    & & \\
cat    & & \\
\end{tabular}
\end{block}
\end{frame}

\begin{frame}[allowframebreaks=0.6]{Activity: SMS Word Count}

\begin{enumerate}
\item Download the file containing sms messages from \url{http://uws.assembla.me/SMSwords.txt} using wget 
\cmd{wget \url{http://uws.assembla.me/SMSwords.txt}}  
\item Sort the file.  \cmd{sort SMSwords.txt} 
\item Sort and remove duplicates.   \cmd{sort SMSwords.txt | uniq} 
\item We need to use a particular option of \lit{uniq} which counts the
number of duplicates, read the manual page for \lit{uniq}. Press \lit{q}
to leave the manual page. \cmd{man uniq} 
\item Sort and count words, \cmd{sort SMSwords.txt | uniq -???} 
\item Sort by the frequency, so that more frequent words appear later, 
\cmd{sort SMSwords.txt | uniq -??? | sort -n}
\end{enumerate}
\end{frame}
%}}}

\section{Summary} %{{{ 

%\begin{frame}
%\ftitle{Summary}
%\end{frame}

Finally, after we are done with what we need to do in the UNIX environment, we
have to logout of the system.  

% ACTIVITY
\begin{frame}{Activity: Logging out of \lit{sunfire}}
To log out of \lit{sunfire}, use the \lit{logout} command, \cmd{logout}
\end{frame}

\begin{frame}{Useful programs/websites}
\begin{itemize}
\item KiTTY, SSH client for Windows\\ \url{http://www.9bis.net/kitty/}
\item Cygwin, UNIX-like environment for Windows\\ \url{http://www.cygwin.com/}
\item Description of computing facilities in SoC\\ \url{http://docs.comp.nus.edu.sg/cf}
\item mySoC, web service portal\\ \url{https://mysoc.nus.edu.sg}
\end{itemize}
\end{frame}
%}}}

\end{document}














\section{Unused content}

\begin{frame}
\ftitle{Acknowledgement}
The materials for this workshop are adapted from the following sources:
\begin{itemize}
\item UNIX Workshop 2005 notes by Mark Tan (SoC, NUS)
\item CS1101 Lab 0 notes by Aaron Tan (SoC, NUS)
\item UNIX/Linux Tutorial for Beginners by Michael Stonebank (University of Surrey)
\end{itemize}
\end{frame}

% Activity to share files 
The following activity, shows you how to share a file on your personal website.  

\begin{frame}[allowframebreaks=0.6]{Activity: \lit{sunfire} as a web host}
\begin{enumerate}
\item Find out your home directory using the following command
\cmd{pwd}
\item Create a directory under your home directory named
\lit{public\_html}. \cmd{mkdir public\_html}
\item Move \lit{hello.txt} into the \lit{public\_html} directory. \cmd{mv hello.txt public\_html}
\item Change the permissions on \lit{public\_html} so that it is readable by
the web server. \cmd{chmod a+rx public\_html}
\item Change your working directory to \lit{public\_html}. \cmd{cd public\_html}
\item Change the permission on \lit{hello.txt} so that it is readable by the
web server.  \cmd{chmod a+r hello.txt}
\item Ask the person next to you to download hello.txt from your personal
website at \url{http://www.comp.nus.edu.sg/~melvin/hello.txt}
\end{enumerate}
\end{frame}


% CLI vs GUI 
% CLI
% + power of piping and redirection, greater control
% + portable (available in many systems)
% + least overhead
% + easy remote access
% + faster to get things done
% + can be easily automated/scripted
% - higher learning curve
% - not easy to multi-task

% GUI
% + display of dynamically modifiable information
% + aesthetics
% + media
% + true multitasking
% + more intuitive, easier to learn

%* You can do things very quickly on command line.
%* GUI prevents you from understanding and learning the functionality happening behind the scenes.
%* Repetitive things can be automated easily using command line.

% GUI: WYSIWYG = WYSIAYG
% CLI: suited to performing ad hoc operations, to quickly combine a few commands
% to perform a query or some tasks

\begin{frame}
\ftitle{CLI vs GUI}
\begin{block}{Graphical User Interface}
What they tell you, WYSIWYG (What you see is what you get). \pause
The truth is, WYSIAYG (What you see is all you get).
\end{block}

\pause

\begin{block}{Command Line Interface}
\begin{itemize}
\item Suited to performing ad hoc operations by combining a few commands.  
\item Easy to automate repetitive tasks.
\item Default interface when accessing remote servers.  
\end{itemize}
\end{block}
\end{frame}



\begin{frame}
\ftitle{Privileges}
Your new UNIX account comes with the following privileges
\begin{itemize}
\item Email : \lit{melvin@comp.nus.edu.sg}
\item Website : \url{http://www.comp.nus.edu.sg/~melvin}
\item Solaris zone: \url{melvin-z.comp.nus.edu.sg}
\item Disk quota: 2Gb
\item Print quota: 50 pages/month
\end{itemize}
\end{frame}

\section{Origins of UNIX}
% people, time, place (story of invention)

\begin{frame}
\ftitle{Creators of UNIX}
\mode<presentation>{
\begin{figure}
\begin{center}\includegraphics[width=0.7\linewidth]{dennis-ken2}\end{center}
\caption{Dennis Ritchie (standing) and Ken Thompson working on a
PDP-11.}
\end{figure}
}
\end{frame}

\begin{frame}
\ftitle{Creators of UNIX}
\begin{figure}
\begin{center}\includegraphics[width=0.7\linewidth]{dennis-ken3}\end{center}
\caption{Ken Thompson (left) and Dennis Ritchie receiving the National Medal of
Technology from President Clinton.}
\label{fig:medal}
\end{figure}
\end{frame}

The first version of UNIX was developed by Ken Thompson in August of 1969,
supposedly in just three week. In 1971, Dennis Ritchie invented the C
programming language and UNIX was rewritten using C and became the first
operating system to be written in a high level language. Most operating systems
of the time, including the original UNIX, was written in assembly language,
which was a special language that is unique to a particular type of hardware.
After UNIX was rewritten in C, it was possible to run UNIX on many different
types of hardware.     

For the development and implementation of UNIX and their contributions to
operating systems theory, Dennis and Ken were jointly awarded the A. M. Turing
award in 1983 and the National Medal of Technology (Figure \ref{fig:medal}) in
1999. The A. M. Turing award is the equivalent of the Nobel prize in Computer
Science.  
% applications
% benefits and uses of UNIX account
\begin{frame}
\ftitle{UNIX in SoC}
The UNIX environment provided by the Solaris OS on our servers are used for:
\begin{itemize}
\item writing programs for your programming labs/assignments
\item learning about operating system concepts (CS2106, Operating Systems)
\item hosting a database driven site (CS2102, Database Systems)
\item accessing SoC printers and checking your print quota
\item reading your SoC email account
\end{itemize}
\end{frame}

% applications
% benefits and uses of UNIX account
\begin{frame}
\ftitle{UNIX in SoC}
The UNIX environment provided by the Solaris OS on our servers are used for:
\begin{itemize}
\item writing programs for your programming labs/assignments
\item learning about operating system concepts (CS2106, Operating Systems)
\item hosting a database driven site (CS2102, Database Systems)
\item accessing SoC printers and checking your print quota
\item reading your SoC email account
\end{itemize}
\end{frame}

\begin{frame}
\ftitle{What is an Operating System?}
\begin{figure}
\begin{center}
\includegraphics[width=0.7\linewidth]{os-layers}
\end{center}
\caption{Relation between applications, OS and hardware}
\label{fig:os}
\end{figure}
\end{frame}

An OS is the software that is in charge of scheduling resources and processes 
of a system as well as managing and maintaining interactivity between the
software and hardware of the system. Figure \ref{fig:os} shows that the
operating system acts as the bridge between your applications and the computer
hardware.  
\section{UNIX Concepts}
UNIX consists of the following components:
\begin{description}
\item[Kernel] The kernel of UNIX is the core of the operating system. It
allocates time and memory to programs and handles the file system and
communications in response to system calls.  
\item[Shell] The shell acts as an interface between the user and the kernel.
When a user logs in, the login program checks the user name and password, and
then starts another program called the shell.   The shell is a command line
interpreter. It interprets the commands the user types in and arranges for them
to be carried out. The commands are themselves programs: when they terminate,
the shell gives the user another prompt.  
\item[Files and processes] Everything in UNIX is either a file or a process.  A
process is an executing program identified by a unique process identifier.  A
file is a collection of data. They are created by users using text editors,
running compilers etc.
\item[Utilities] An operating system is useless without useful applications that
help the you to accomplish everyday tasks such as editing text, browsing the web
etc.  A UNIX-like environment provides a standard set of useful utilities.  
\end{description}

\section{Checking your SoC email account}
\begin{frame}
\ftitle{Checking UNIX email}
% Checking UNIX email, show webmail version, refer to CF site for details on
% setting up mail client, email forwarding from MySoC
You can access your UNIX email account via mySoC Webmail, \url{http://mysoc.nus.edu.sg/~webmail}
\mode<presentation>{
\begin{figure}
\begin{center}
\includegraphics[scale=0.15]{mysoc_webmail}
\caption{mySoC Webmail interface}
\end{center}
\end{figure}
}
Your mailbox part of your disk usage, which is 2Gb. You can forward your NUSNET
email to your UNIX email using \url{https://exchange.nus.edu.sg/autoforward}.
\end{frame}

% vim history
% vim is a modal editor, command mode and insert mode
In this workshop, we will introduce you to the Vim editor.  Vim stands for Vi
IMproved, it is created by Bram Moolenaar (now working at Google) as an
improvement of an earlier editor called vi by Bill Joy (co-founder of Sun
Microsystems).  Vim is a modal editor. This means, depending on the current
mode, different keys have different effects (see Figure \ref{fig:vim_modes}).  
For example in Normal mode, pressing `x' key deletes the character under the
cursor, while in Insert mode, it adds the character `x' to the text file. This
allows Vim to use almost all the keys on your keyboard for issuing commands.  

\begin{frame}
\ftitle{Vim is a modal editor}
\begin{figure}
\begin{center}
\includegraphics[scale=0.4]{vim_modes}
\end{center}
\caption{Different modes of Vim and how to switch between them}
\label{fig:vim_modes}
\end{figure}
\end{frame}


\section{Some useful applications on UNIX}
A major part of the operating system that the user interacts with are the
utilities that helps the user to perform everyday tasks such as editing files,
printing and reading email.     

You have already used a number of such utilities such as \lit{mkdir} (make
directory), \lit{mv} (move), \lit{chmod} (change modifiers/permissions)
and \lit{cd} (change directory).

In the next activity, you will learn a number of utilities related to printing.
There are a number of printers that are exclusive to SoC students, they are
located in COM1 level 1 and basement.  

% ACTIVITY: Printing
\begin{frame}{Activity: Printing}
\mode<presentation>{
\begin{figure}
\begin{center}
\includegraphics[scale=0.4]{printer}
\end{center}
\caption{Lexmark printers at COM1}
\end{figure}
}

\begin{itemize}
\item View the status of the print queue, use \lit{lpq}, \cmd{lpq -P pstsc}
\item Remove a print job after it has been sent, use \lit{lprm}, \cmd{lprm -P pstsc 89}
\item Check your print quota, use \lit{pusage}, \cmd{pusage}
\end{itemize}

\end{frame}

If there is additional time during the workshop, you may wish to setup a
homepage which is accessible at \url{http://www.comp.nus.edu.sg/~melvin}.
The details are provided in the following activity.  

% ACTIVITY: setting up index.html
\begin{frame}[fragile]
\frametitle{Activity: Setting up your homepage (Optional)}
Instead of using \lit{hello.txt}, create a file named \lit{index.html} and
put it in your \lit{public\_html} directory. Remember to change its
permissions to readable by all.  

\begin{verbatim}
<html>
<head>
<title>Sample index page</title>
</head>
<body>
Hello World
</body>
</html>
\end{verbatim}
\end{frame}

\section{Organising your home directory}

\begin{frame}{Command line interface}
\begin{itemize}
\item Easy to automate repetitive tasks
\item Default interface when accessing remote servers %not that important..
\end{itemize}
\end{frame}

\begin{frame}[fragile]
\frametitle{Activity: playing with \lit{diff} and \lit{grep} II}
Let's see how \lit{text} and \lit{text\~} look like first

\begin{columns}
\begin{column}{0.5\textwidth}
\\ Output of \lit{cat text}

\begin{verbatim}
same text
same text
Hello World!
still the same
still the same
\end{verbatim}
\end{column}

\begin{column}{0.5\textwidth}
\\ Output of \lit{cat text\~}
\begin{verbatim}
same text
same text
Hello World~
still the same
still the same
\end{verbatim}
\end{column}
\end{columns}
\end{frame}


\begin{frame}[fragile]
\frametitle{Activity: playing with \lit{diff} and \lit{grep} III}
Output of \lit{diff text text\~} :
\begin{verbatim}
3c3
< Hello World!
---
> Hello World~
\end{verbatim}
\end{frame}


\begin{frame}
\frametitle{Activity: playing with \lit{diff} and \lit{grep} IV}
\begin{enumerate}
\item What is \lit{grep}? *-- look for a pattern in file(s)
\cmd{grep pattern file}
\item Sometimes it is useful to find the occurrences of some word in a (list of) file.\\ 
Say you suspect a typo in you source code,\\ 
open a text editor and 'Find'?
\item But what if you made the same typo in a lot of files?\\
'grep' makes your life easier\\
Let's find out how to 'grep'
\end{enumerate}
\end{frame}


\begin{frame}[fragile]
\frametitle{Activity: playing with \lit{diff} and \lit{grep} V}
Sample output of \lit{grep h1 a.html}:

\begin{columns}
\begin{column}{0.5\textwidth}
\lit{grep h1 a.html}

\begin{verbatim}
    <h1>Hello World!</h1>
\end{verbatim}
\end{column}

\begin{column}{0.5\textwidth}
\\Contents of a.html
\begin{verbatim}
<html>
  <body>
    <h1>Hello World!</h1>
  </body>
</html>
\end{verbatim}
\end{column}
\end{columns}

\pause
\textcolor{red}{A bit too easy, isn't it? Ready to get nasty?}
\end{frame}

\begin{frame}[fragile]
\ftitle{Getting serious about grep}
\begin{enumerate}

\item Output of \lit{grep h1 a.html*}

\begin{verbatim}
a.html:<h1>Hello World!</h1>
a.html~:<h1>Hello World~</h1>
\end{verbatim}

\pause
\item Output of \lit{grep -n h1 a.html*}
\begin{verbatim}
a.html:3:<h1>Hello World!</h1>
a.html~:3:<h1>Hello World~</h1>
\end{verbatim}

\pause
\item Output of \lit{grep -n -i 'heLlO wORlD' a.html*}
\begin{verbatim}
a.html:3:<h1>Hello World!</h1>
a.html~:3:<h1>Hello World~</h1>
\end{verbatim}

\pause
\item Find out more in \lit{'man grep'} !
\pause
\item grep on Linux is more fun! :p
\end{enumerate}
\end{frame}
